\documentclass[11pt]{article}
\usepackage{cocina}
\begin{document}
\begin{ficha}{Comparar}{¿Quién se lleva más?}{2}

\begin{encargo}
Dos mesas, dos pedidos: ¿quién se lleva más pizza? \textbf{Antes de cortar nada,
apuesta} en la casilla: $>$, $<$ o $=$. Luego compruébalo con las pizzas.
\end{encargo}

\bloque{1}{La escalera}{Con pizza}\quad
{\footnotesize apuesta, colorea los pedidos y di por qué}

\vspace{2mm}
{\small
\begin{tabular}{@{}l@{\hspace{3mm}}c@{\hspace{1mm}}c@{\hspace{2mm}}c@{\hspace{2mm}}c@{\hspace{1mm}}l@{\hspace{4mm}}l@{}}
\ej{a} & \fA{3}{8} \pizza[escala=.6]{8}{} & \signo & \pizza[escala=.6]{8}{} \fB{5}{8} && mismos trozos: \hueco[4.4cm]\\[3mm]
\ej{b} & \fA{1}{3} \pizza[escala=.6]{3}{} & \signo & \pizza[escala=.6]{4}{} \fB{1}{4} && un trozo cada uno: \hueco[3.8cm]\\[3mm]
\ej{c} & \fA{3}{4} \pizza[escala=.6,guias=8]{4}{} & \signo & \pizza[escala=.6]{8}{} \fB{5}{8} && $\dfrac{3}{4} = \fhueco[][8]$\quad corta \hueco[.8cm] caja\\[3mm]
\ej{d} & \fA{2}{4} \pizza[escala=.6,guias=12]{4}{} & \signo & \pizza[escala=.6,guias=12]{6}{} \fB{3}{6} && $\dfrac{2}{4} = \fhueco[][12]$\quad $\dfrac{3}{6} = \fhueco[][12]$\\
\end{tabular}}

\vspace{1.5mm}
\begin{nick}
\textbf{Con trozos iguales}, basta contar trozos. \textbf{Con el mismo número de
trozos}, gana el trozo más grande: el de la pizza cortada en menos partes. \textbf{Con
trozos distintos}, contar no vale: se cortan las dos en la \textbf{misma secuencia
de trozos} —eso es el \textbf{denominador común}— y entonces sí, se cuenta.
\end{nick}

\bloque{2}{Ninguno cabe en el otro}{Con la regla}\quad
{\footnotesize ni los tercios caben en los quintos ni al revés}

\vspace{1mm}
{\small
\begin{minipage}[c]{.42\linewidth}
\cortadohueco{2}{3}{5}\qquad\cortadohueco{3}{5}{3}
\end{minipage}%
\begin{minipage}[c]{.58\linewidth}
Mesa 4: \fA{2}{3}\quad Mesa 7: \fB{3}{5}\quad\ \fA{2}{3}\ \signo\ \fB{3}{5}\\[1.5mm]
Márcalas en la regla:\\[1mm]\reglavacia[8.6cm]{1}{15}\\[1mm]
Gana por \hueco[1.4cm] de pizza.
\end{minipage}}

\bloque{3}{Sin pizza}{Sobre el papel}\quad
{\footnotesize busca la secuencia común, cuenta y pon el signo}

\vspace{1mm}
{\small\renewcommand{\arraystretch}{2.4}
\begin{tabular}{@{}l@{\hspace{10mm}}l@{}}
\ej{a}$\dfrac{5}{6}\ \signo\ \dfrac{7}{9}$\qquad en \hueco[.8cm]:\quad $\fhueco[][]\ \ \fhueco[][]$ &
\ej{b}$\dfrac{3}{10}\ \signo\ \dfrac{1}{4}$\qquad en \hueco[.8cm]:\quad $\fhueco[][]\ \ \fhueco[][]$\\
\ej{c}$\dfrac{4}{9}\ \signo\ \dfrac{5}{12}$\qquad en \hueco[.8cm]:\quad $\fhueco[][]\ \ \fhueco[][]$ &
\ej{d}$\dfrac{7}{12}\ \signo\ \dfrac{5}{8}$\qquad en \hueco[.8cm]:\quad $\fhueco[][]\ \ \fhueco[][]$\\
\end{tabular}}

\alto{el número de abajo dice \textbf{en cuántas partes} se corta la pizza, no
cuánto te toca. $\frac{1}{4}$ lleva un 4 y $\frac{1}{3}$ un 3, pero el cuarto es
más pequeño: más comensales, menos ración. Y contar trozos que no miden lo mismo
no sirve para nada.}

\reto{¿Hay alguna fracción entre $\frac{3}{4}$ y $\frac{4}{5}$? Encuentra una.
¿Y entre $\frac{1}{2}$ y $\frac{2}{4}$?}

\comprueba{Cada código abre ese pedido en Practicar.}{%
\qr{1a}{j=comparar\&a=3/8\&b=5/8}\hfill\qr{1b}{j=comparar\&a=1/3\&b=1/4}\hfill
\qr{1c}{j=comparar\&a=3/4\&b=5/8}\hfill\qr{1d}{j=comparar\&a=2/4\&b=3/6}\hfill
\qr{2}{j=comparar\&a=2/3\&b=3/5}}

\end{ficha}

% ─────────────────────────────── REVERSO ──────────────────────────────────
\begin{dorso}{Comparar}{Más mesas}{2}

\bloque{1}{En fila}{Con la regla}\quad
{\footnotesize márcalas en la regla y ordénalas de menor a mayor}

\vspace{1mm}
{\small
\ej{a}$\dfrac{1}{2}$,\ $\dfrac{3}{8}$,\ $\dfrac{3}{4}$,\ $\dfrac{5}{8}$\qquad
\reglavacia[10.5cm]{1}{8}\\[1mm]
\hspace*{6mm}\hueco[.9cm] $<$ \hueco[.9cm] $<$ \hueco[.9cm] $<$ \hueco[.9cm]

\vspace{4mm}
\ej{b}$\dfrac{2}{3}$,\ $\dfrac{1}{2}$,\ $\dfrac{5}{6}$,\ $\dfrac{1}{3}$\qquad
\reglavacia[10.5cm]{1}{6}\\[1mm]
\hspace*{6mm}\hueco[.9cm] $<$ \hueco[.9cm] $<$ \hueco[.9cm] $<$ \hueco[.9cm]}

\bloque{2}{¿Qué peldaño?}{Sobre el papel}\quad
{\footnotesize di qué peldaño de la escalera toca, y pon el signo}

\vspace{1mm}
{\footnotesize\color{tenue} A: mismos trozos \quad B: mismo número de trozos \quad
C: uno cabe en el otro \quad D: ninguno cabe \quad E: iguales disfrazados}

\vspace{1mm}
{\small\renewcommand{\arraystretch}{2.8}
\begin{tabular}{@{}l@{\hspace{6mm}}l@{\hspace{14mm}}l@{\hspace{6mm}}l@{}}
\ej{a}$\dfrac{4}{7}\ \signo\ \dfrac{4}{9}$ & peldaño \hueco[.7cm] &
\ej{b}$\dfrac{5}{12}\ \signo\ \dfrac{7}{12}$ & peldaño \hueco[.7cm]\\
\ej{c}$\dfrac{2}{5}\ \signo\ \dfrac{7}{10}$ & peldaño \hueco[.7cm] &
\ej{d}$\dfrac{3}{4}\ \signo\ \dfrac{5}{6}$ & peldaño \hueco[.7cm]\\
\ej{e}$\dfrac{6}{9}\ \signo\ \dfrac{4}{6}$ & peldaño \hueco[.7cm] &
\ej{f}$\dfrac{9}{10}\ \signo\ \dfrac{9}{8}$ & peldaño \hueco[.7cm]\\
\end{tabular}}

\bloque{3}{Problemas de las mesas}{Sobre el papel}

\vspace{1.5mm}
{\small
\ej{a}En la mesa 4 se comen $\frac{3}{5}$ de una pizza y en la mesa 7,
$\frac{5}{8}$ de otra igual. ¿Quién ha comido más? \hueco[2cm]\quad ¿Por cuánto?
\hueco[1.4cm]\par\vspace{12mm}

\ej{b}Ana ha leído $\frac{2}{3}$ de un libro y Luis $\frac{5}{8}$ del mismo
libro. ¿Quién va más adelantado? \hueco[2cm]\par\vspace{12mm}

\ej{c}En una clase aprueban $\frac{3}{4}$ de los alumnos; en otra,
$\frac{7}{10}$. ¿En cuál aprueba una parte mayor de la clase? \hueco[2cm]\par\vspace{12mm}}

\reto{Nick dice: «de dos fracciones, es más pequeña la que tiene más distancia
entre el número de arriba y el de abajo». Prueba con $\frac{2}{3}$ y
$\frac{3}{4}$. ¿Funciona siempre? Busca un caso en que falle.}

\comprueba{Los del bloque 2 abren la pareja; los del 3, la llamada al teléfono, porque el cortador no llega.}{%
\qr{2b}{j=comparar\&a=5/12\&b=7/12}\hfill\qr{2c}{j=comparar\&a=2/5\&b=7/10}\hfill
\qr{2d}{j=comparar\&a=3/4\&b=5/6}\hfill\qr{3a}{j=telefono\&a=3/5\&b=5/8}\hfill
\qr{3b}{j=telefono\&a=2/3\&b=5/8}\hfill\qr{3c}{j=telefono\&a=3/4\&b=7/10}}
\end{dorso}
\end{document}
